\documentclass[11pt]{article}

% ============================================================
% Explanation of Revisions (ARR resubmission, compulsory upload)
%
% Build:  latexmk -pdf EXPLANATION_OF_REVISIONS.tex
%
% TWO SWITCHES. Both must be off for the posted PDF.
%
%   \draftnotestrue   shows internal decisions and the pre-submission
%                     checklist in red. For our eyes only.
%   \draftnotesfalse  hides them. THIS IS THE POSTING BUILD.
%
% Locations are given as section and appendix numbers only. No line
% references, so this document does not depend on a frozen main.pdf and
% survives recompilation of main.tex. The \ph{} macro is retained but
% unused. Do not reintroduce line numbers.
%
% ANONYMITY: no authors, no affiliations, no repository URLs, and no
% author-list change information. Keep it that way.
% ============================================================

\newif\ifdraftnotes
\draftnotesfalse       % <<<< POSTING BUILD. Set \draftnotestrue to see internal notes.

\usepackage[T1]{fontenc}
\usepackage[utf8]{inputenc}
\usepackage{times}
\usepackage{microtype}
\usepackage[margin=1in]{geometry}
\usepackage{booktabs}
\usepackage{longtable}
\usepackage{array}
\usepackage{enumitem}
\usepackage{xcolor}
\usepackage{amsmath}
\usepackage{amssymb}
\usepackage[hidelinks]{hyperref}

\setlist[itemize]{leftmargin=1.2em,itemsep=0.35em,topsep=0.35em}
\setlength{\parskip}{0.4em}
\setlength{\parindent}{0pt}

% Internal note. Hidden in the posting build.
\newcommand{\note}[1]{%
  \ifdraftnotes\par\smallskip{\color{red}\small\textbf{[}#1\textbf{]}}\par\smallskip\fi}
% Unresolved placeholder. Highlighted in the draft build.
\newcommand{\ph}[1]{\ifdraftnotes{\color{red}#1}\else#1\fi}

\title{\vspace{-2.5em}\textbf{Explanation of Revisions}\vspace{-0.8em}}
\date{}

\begin{document}
\maketitle

Resubmission of an ARR 2026 May submission.

This is a substantial revision. The entire discussion-phase campaign is now
in the manuscript, nine experiment families across the model suite, covering
direct-geometry baselines, intervention validation, cross-seed stability,
grouping recovery, a blinded label audit, matched-KL editing controls, error
tails, a lower sparsity budget, and the cross-lens study. \S3 of the paper
was rebuilt as a single argument and \S4 is new.

\S4 carries the finding. Two Jacobian lenses differing only in fitting
corpus, one English and one Chinese, read the same frozen Qwen3.5-9B states
and report different scripts on 9/10 factual-recall prompts, while one SRP
feature stays dominant under both readings on \textbf{77/80 prompts (96\%,
95\% CI [0.90, 0.99])}. What a fitted lens reports follows the corpus it
was fitted on, so a token reading alone leaves the working language of a
computation underdetermined. The decomposition is what makes two competing
readings comparable in the first place, and it supplies the control that
settles them. This is the systematic new interpretation the meta-review
asked for, and \S4 applies it to a reading published in contemporaneous work.

Alongside it, the revision adds the downstream-reliability evidence the
reviews asked for.
Predicted against realised margin change reaches $r^2=0.83$--$0.93$ on six
models, with random-direction controls at most 0.06 (\S3.1, Table~2).

Reviewers returning to the paper will find the argument in \S1, \S3 and \S4.
\S2 is extended but unchanged in substance.

This document describes what changed and where. Every number below appears in
the revised manuscript at the section or appendix given.

\note{DRAFT BUILD. Set \texttt{\textbackslash draftnotesfalse} and strip
every red note before converting the posting PDF.}

\section{Overview}

\begin{itemize}
\item The previous round asked for three things. Evidence of what the method
enables \emph{systematically} beyond token rankings, the discussion-phase
experiments moved into the manuscript so they can be reviewed, and the core
objects defined early. The paper is restructured around all three.
\item Every experiment reported during the discussion phase now has a named
manuscript location in \S3, \S4, or Appendices D--N, listed in the map
below.
\item Responsible NLP Checklist items C1, on compute budget and hardware, and
C4, on package versions, are now answered in the paper, and the code and
trained dictionaries are released.
\item All quantitative claims were re-verified against the underlying result
artifacts.
\end{itemize}

\section{Response to the meta-review}

The meta-review asks what the method brings to the interpretability field.
The revision answers by argument as well as by experiment. The paper is now
built on one claim, that token identity is too coarse a unit for interpreting
vocabulary readouts, and \S1 closes by stating two contributions, the
interface and the cross-lens finding. Sections~3 and~4 are organised to test that claim
rather than to demonstrate a tool.

\subsection*{Request 1. ``What new interpretations can be done in a
systematic way?'', for instance features aligning with different senses}

\begin{itemize}
\item \textbf{The sense-alignment test the meta-review proposed is run
directly, and it passes.} On CoarseWSD-20, feature groups selected on
training contexts predict held-out senses at balanced accuracy \textbf{0.90 /
0.92 / 0.80} across three models, against a label-permutation null of 0.41
(\S3.4, Appendix~H). Scoped as an alignment check rather than a
disambiguation-benchmark claim, it shows that the discovered features carry
sense structure and not only token form. Appendix~H reports the
comparison probes.
\item New top-level \S4, \emph{Corpus Conditionality}. Two Jacobian lenses
differing only in fitting corpus, 100
English prompts against 100 Chinese, read the same frozen Qwen3.5-9B states.
They report different scripts on 9/10 factual-recall prompts, while one SRP
feature stays dominant under both readings on \textbf{77/80 prompts (96\%,
95\% CI [0.90, 0.99])}. The null floors are 0/159 for within-lens unrelated
tokens and 0/240 for shuffled pairings. Appendix~K has the protocol and the
per-family tables.
\item Three extensions establish that the effect is not specific to one
construction or corpus size (\S4.3). A same-script
English--German pair with cognates excluded, a second lens family of
per-layer ridge translators in the tuned-lens style, and a tripled fitting
corpus.
\item \S4 applies this control to a multilingual reading reported in
contemporaneous work from July 2026, under the same lens construction and the
same prompt, and finds the reported language moves with the fitting corpus
(\S4.4). The related-work section and Appendix~P position the
paper against that work. The section and the citation are new in this
revision.
\item The worked case the discussion thread centred on is written up in full.
Appendix~K gives the published antonym prompt layer by
layer. The shared feature contributes $+12.5$ and $+13.4$ logits at layer 24
under the English-fitted and Chinese-fitted lenses, and $+14.5$ and $+17.6$
at layer 29.
\end{itemize}

\subsection*{Request 2. Discussion-phase experiments ``should be included in
a new revision''}

Every experiment reported during the discussion phase is now in the
manuscript, in the main text or in a numbered appendix, with its protocol,
its controls and its per-model tables. None is carried by summary alone and
none was dropped. The map below names the location of each, so every
one of them can be checked against the paper rather than against the response
thread.

\subsection*{Request 3. Define contrast and readout early}

\emph{Readout}, \emph{selected readout score}, \emph{contrast}, and
\emph{margin} are all defined in \S1, ahead of the method and every result, and
used consistently thereafter.

\section{Reproducibility and released artifacts}

\begin{itemize}
\item \textbf{Compute and hardware (checklist C1).} Appendix~F gives the
hardware, a single NVIDIA A40 per job on an academic Slurm cluster with a
24\,GB-class GPU sufficient for evaluation. It also gives per-model training
cost of 6--12 A40 GPU-hours per readout SAE, run once per model, the
evaluation cost, and the local-inference configuration.
\item \textbf{Package settings (checklist C4).} Appendix~F pins the Python
version at 3.11--3.12, pins the package versions, and points at the lockfile
in the released code, which also carries per-model training and evaluation
recipes.
\item \textbf{Artifacts.} The revision releases, under an MIT license, the
code with per-model training and evaluation recipes, the selected readout SAE
checkpoints for all eight models at both the $k{=}128$ and $k{=}256$
operating points, the contrast banks, both the diagnostic bank and the
benchmark-derived bank, and the blinded label-audit protocol (Appendix~F).
\end{itemize}

\section{Where each discussion-phase experiment now lives}
\label{sec:map}

\begingroup
\small
\setlength{\tabcolsep}{4pt}
\renewcommand{\arraystretch}{1.15}
\begin{longtable}{@{}p{0.24\textwidth}p{0.475\textwidth}p{0.215\textwidth}@{}}
\toprule
\textbf{Experiment} & \textbf{Result} & \textbf{Location} \\
\midrule
\endfirsthead
\toprule
\textbf{Experiment} & \textbf{Result} & \textbf{Location} \\
\midrule
\endhead
\bottomrule
\endfoot

Cross-lens study
& Two lenses differing only in fitting corpus report different scripts on
9/10 prompts, while one feature stays dominant under both readings on 77/80
(96\%, 95\% CI [0.90, 0.99]). Null floors 0/159 and 0/240
& \S4. App.~K \\

Intervention validation, six models
& Predicted against realised margin change at $r^2=0.83$--$0.93$,
random-direction controls at most 0.06, slopes reported as measured
& \S3.1, Table~2. App.~J \\

Direct-geometry baselines (k-means at two widths, hard clustering, row-kNN,
nearest-row ridge, PCA-256), six models
& SRP coverage 0.72--0.81 against at most 0.67, leading the strongest
per-model alternative by 8.9--17.3 points, intervals non-overlapping
& \S3.2, Table~1. Full method grid in App.~G \\

Held-out grouping recovery
& Same-recipe held-out dictionary recovers 72--75\% of the seed-stable
grouping core, against 43--49\% for row-kNN and 21--25\% for clustering
& \S3.3. App.~I.3 \\

Cross-seed stability
& Every tested contrast reproduces above the null's p90 at both widths, at
roughly 15 times the null. 87--91\% of used feature groups have an
above-chance counterpart
& \S3.3. App.~I \\

Matched-KL editing comparison
& SRP beats mean-row, rank-4 PCA and PCA group-direction controls at all four
matched-KL targets on the display model, 12 of 12 paired with all CIs
excluding zero. Per-model results reported
& \S3.4. App.~N \\

Blinded label audit
& Three runs, frozen prompt, top-20 rows only. 93.4\% unanimity, Fleiss'
$\kappa=0.87$, and 81.8\% agreement with the author audit
& \S3.4. App.~L \\

Error tails and margin stratification
& Mean, p95 and maximum absolute error for SRP and all six alternatives, plus
a margin-stratified table
& App.~G.4 \\

Lower sparsity budget, two models at $k=64$ and $k=32$
& A smaller budget strands capacity rather than buying sparsity. At matched
width, $k=64$ leaves 32--51\% of features dead against under 4\% at $k=256$
and $k=32$ leaves 64--69\%, with rowEV and KL degrading on both models
& App.~D.3 \\
\end{longtable}
\endgroup

The restructuring moved the appendix letters, so the pointers given in the
discussion responses no longer resolve. Among others, the benchmark-derived
and prompt-injection decompositions cited as App.~I.4 and I.7 are now in
Appendix~L, the constrained edit test cited as App.~K is now Appendix~N, the
diagnostics cited as App.~G.1 are now Appendix~F, and the decomposition
identities cited as App.~B.2 are now Appendix~A. The appendix roadmap on the
first appendix page maps all sixteen.

\section{Point-by-point}

Old line numbers refer to the previous submission, as cited in the reviews.

\subsection*{Reviewer j4yY}

\begin{itemize}
\item \textbf{Abstract, old L24. Why do logit differences matter, and what is
the end goal?} The Introduction now states the objective in its first two
paragraphs. A logit difference is the model's exact signed
preference between two candidate outputs, and the paper's target is what
carries that preference on the readout side. The abstract was rewritten and
now opens on the ambiguity of the token as a unit, so the motivation for
reading logit differences arrives before the method does. The fidelity the
paper rests on is sign agreement of at least 0.91 suite-wide and
sign-preserved low-error coverage of 0.72--0.81 on the six gated models
(\S3).

\item \textbf{Intro, old L84. Does the basis work across models, questions,
and domains?} One basis per model serves all prompts, contrasts and lenses
with no refitting, and the same recipe applies unchanged to all eight LM
heads. The 300 benchmark-derived contrasts across QA, contract-entailment and
security formats reach sign agreement 0.94 and coverage 0.84 on Qwen3.5-2B,
and 0.98 and 0.93 on Qwen3.5-0.8B (\S3 Setup, App.~L). Table~1 now
makes the 8-model reconstruction suite and the 6-model interpretation gate
explicit.

\item \textbf{Intro. Define ``contrast'' and ``readout'' early.} Both are
defined in \S1, along with \emph{selected readout score} and
\emph{margin}.

\item \textbf{Method. Why not absorb the shared offset $\mu$?} \S2 defines $\mu$ and shows its exact cancellation whenever the
coefficients sum to zero, which covers every contrast used, and Appendix~A
derives the identity. A new paragraph in Appendix~A.1, \emph{Why the offset
stays outside the dictionary}, answers the question directly.
Folding $\mu$ into the dictionary would give an atom carrying coefficient one
on every row, so it would occupy a TopK slot for every token and spend sparse
capacity on a quantity that is identical everywhere, and its decoder
direction would be drawn toward the row mean. Keeping the offset explicit
follows standard sparse autoencoder practice and the documented stabilising
effect of centering output embeddings.

\item \textbf{Method. Is the residual orthogonal to the dictionary
directions?} It is not, and with an overcomplete dictionary spanning the
space it cannot be. \S2 now describes the separation as additive
and exact, alongside the reading rule that gates interpretation on residual
size. Appendix~A gives the identity. Appendix~F quantifies what orthogonality
would buy, since least-squares refitting on the encoder's fixed support
raises rowEV by only 0.004--0.006, so the non-orthogonal residual is not
absorbing recoverable structure.

\item \textbf{Method, old L196. How are contrasts justified, and how do we
know A is better than B?} \S3 Setup states plainly that a
positive margin says the readout favours A over B and that correctness is a
separate question, and that only the separate benchmark-derived bank takes A
from a gold label. The banks are fixed once, before any decomposition is
read. Appendix~G covers bank families, provenance and filters.

\item \textbf{Experiment, old L302. Is sign agreement of 0.91 sufficient, and
what about decision boundaries and generation?} The reading rule (\S2) gates interpretation on sign agreement and relative error, and
only accepted scores are read qualitatively. Sign errors concentrate at
near-tie margins (App.~G.4). New in this revision,
the intervention study of Table~2 supplies the downstream-reliability
evidence the question asked for.

\item \textbf{Table 3. Why median rather than mean, and what about tail
failures?} Appendix~G.4 is new and reports mean, p95 and maximum absolute
error for SRP and all six alternatives, plus a margin-stratified table. SRP
has the lowest value on every statistic on every model, and the remaining
tail sits at near-indifference margins. \S3.2 carries a
one-sentence pointer.

\item \textbf{Question 3 and summary. Novelty relative to an SAE.} The
contribution split is stated in \S1. TopK is the fitting
primitive, and the contribution is the factorised readout object and the
interface it induces. The Introduction separates SRP from activation SAEs,
residual-stream attribution and direct LM-head geometry by what each
decomposes. Table~1 and \S3.3 give the
empirical case.

\item \textbf{Summary. No downstream-task analysis.} Three responses. The 300
benchmark-derived contrasts above, the intervention study of Table~2, and
\S4, which changes how a published lens reading should be used.
\end{itemize}

\subsection*{Reviewer e1hF}

\begin{itemize}
\item \textbf{W1. Does the learned basis beat direct analysis of LM-head
rows?} We ran the exact comparison requested, including both constructions
the review named, clustering the rows of $W_U$ and treating top-k nearest
neighbours as a basis. Table~1 gives six direct-geometry
alternatives on all six gated models, among them a capacity-matched k-means
dictionary at SRP's own $D$ and $k$. \S3.3 tests the review's
stated criterion directly, whether the resulting groupings look similar to
SRP's. Scored on identical cores, row-kNN recovers 43--49\% and clustering
21--25\%, against 72--75\% for a held-out same-recipe SRP dictionary. The
antecedent does not hold.

\item \textbf{W2. Are the explanations stable across SAE runs?} Three seeds
at each of two widths, with equal held-out reconstruction verified first, so
the review's premise of a different seed with equally good reconstruction is
controlled and not merely assumed (App.~I). What reproduces is the object the
review cares about, the decomposition of a selected score, with 100\% of
contrasts above the null's p90 at roughly 15 times the null and negligible
cross-side leakage, together with the token-group structure, where 87--91\%
sit above the null's p99. Individual decoder directions do not recur, at a
matched cosine of roughly 0.33, which is the expected signature of feature
splitting and is stated in the revision. Claims are scoped within-recipe, and
Appendix~I.4 reports the cross-recipe comparison.

\item \textbf{W3. Simpler controls, and matched KL rather than fixed scale.}
We added mean-row, rank-4 PCA and PCA group-direction controls, all built
from the same five discovery rows without touching held-out tokens, then
swept and interpolated them to four common KL targets (App.~N). SRP wins all
12 paired comparisons on the display model, with every 95\% bootstrap
interval excluding zero. Against the strongest control, the mean-row
direction, the paired advantage runs from $+0.032$ to $+0.087$ across the
four KL targets. On the review's specific objection, the KL gap of 0.679
against 0.004 for the oracle, Appendix~N now explains why that gap does not
measure efficiency. The oracle bias penalises the logits of the evaluation
lexicon itself, so it moves exactly the tokens being scored and leaves the
rest of the distribution untouched, while a readout direction selected from
five discovery tokens acts through the whole vocabulary and pays for whatever
else it moves. The oracle is an upper bound that reads the evaluation labels,
and the comparisons carrying the claim are the label-free mean-row and PCA
directions, which are matched on distributional cost. We state the boundary
in the main text, and \S3.4 scopes the advantage to the display
model, with the per-model results in Appendix~N.

\item \textbf{Comment. Is $k=128/256$ sparse enough, given that activation
SAEs use 32 or 64?} Appendix~D.3 is new and retrains Qwen3.5-2B and
Qwen3.5-0.8B at both values the comment names, $k=64$ and $k=32$, under an
otherwise identical recipe. Going from $k=256$ to $k=64$ raises the
dead-feature rate from under 4\% to 32--51\%, and $k=32$ strands 64--69\% of
the dictionary, with rowEV falling from 0.80 to 0.68 and KL rising from 0.23
to 0.37 on the display model. A smaller budget strands capacity instead of
buying sparsity, and the progression repeats on both models. Appendix~D adds the
fraction-active framing, since absolute $k$ is not comparable across
dictionaries of different width. At $D=65{,}536$, $k=128$ and $k=256$
activate 0.20\% and 0.39\% of the dictionary, and a median of 65--107
features carries 80\% of a score's contribution mass.
\end{itemize}

\subsection*{Reviewer pZ3j}

\begin{itemize}
\item \textbf{W1. The novelty of the method seems limited.} Same treatment as
j4yY's Question 3. The contribution split is in \S1, the boundary against
activation SAEs and dense decompositions is in the Introduction,
and the distinctness evidence is in Table~1 and \S3.3.

\item \textbf{W2. Labelling is human, many labels are ambiguous, and the
success rate is unknown.} Labelling is not done by hand. Appendix~L now
describes the pipeline, where a language model drafts each description from
the feature's highest-scoring unembedding rows and the authors verify every
label used to support a claim against those same rows. The success rate is
the rate at which that verification succeeds, and it is now measured on two
axes. On intrinsic labelability, 85 of 99
claim-bearing display labels name a coherent lexical family at top-20 row
depth. On panel relevance, 66 of those 85 are on-target for the sense their
panel claims, all side-consistent, with no case where an on-target label
opposed the prose. Both appear in Appendix~L. A blinded
three-run audit with a frozen prompt, fresh contexts, and no panel claims or
author labels in view reaches 93.4\% unanimity, $\kappa=0.87$, and 81.8\%
agreement with the author audit, with disagreements skewing toward
\emph{more} coherent. We report it as a same-model reproducibility audit,
since three runs of one model do not constitute three annotators. \S3.4 states that labels are readability summaries kept separate from
the measured objects, which are feature ids, signed terms and residuals, so a
decomposition does not depend on accepting a label.

\item \textbf{W3. Many hyperparameters, and the sensitivity is not analysed.}
Two widths and three seeds run throughout the stability analyses, with sweeps
over width, sparsity, architecture, initialisation and schedule in
Appendices~C--E and Appendix~D.3 new. The main text surfaces this at \S2
(operating-point rule), which describes the reported basis as an
audited operating point and disclaims canonical status.

\item \textbf{W4. No genuinely new analysis that required the tool.} \S4 is
that analysis. It re-tests a reading published in contemporaneous work, finds
the reported language moves with the lens fitting corpus, and identifies the
one feature that survives both readings on 77/80 prompts. The comparison
requires a basis shared across the two lenses, which is what the
decomposition supplies. See the response to the meta-review above.
\end{itemize}

\section{Additional changes}

\begin{itemize}
\item \textbf{Section 3 rebuilt as one argument.} It now runs from fidelity
to distinctness against direct-geometry alternatives, then to reproducibility
and intervention relevance, with the evaluation design folded into its
opening.
\item \textbf{Supporting material moved to the appendices.} Feature-resolved
DLA, the lexical-edit protocol, the score-family grid and additional
qualitative displays are now in Appendices~M, N, G and L, making room in the
main text for \S4 and the new evaluations. The main text keeps a
one-paragraph summary of each with scoped claims.
\item \textbf{Scope statements made explicit.} Edit claims report per-model
results, intervention claims are worded as local readout-side predictions,
and the label audit is described as a blinded same-model reproducibility
audit rather than human annotation.
\item \textbf{Notation.} Floored and unfloored relative error now carry
distinct symbols.
\end{itemize}

The revision moves the paper from a decomposition with diagnostics to a claim
with evidence. The interface is now tested against direct-geometry
alternatives, against seeds, and against intervention, and \S4 uses it to
show that what a published lens reports follows the corpus the lens was
fitted on. That last point is a fact about fitted lenses rather than
about this method, and reaching it took a basis the two lenses could be read
in.

% ============================================================
% Everything below is internal and never appears in the posting build.
% ============================================================
\ifdraftnotes
\clearpage
{\color{red}
\section*{Pre-submission checklist (internal, not for posting)}

\textbf{Manuscript}
\begin{itemize}[label=$\square$]
\item[$\boxtimes$] Abstract. \textbf{Decided 2026-08-04, leave as is.} The
abstract keeps row explained variance of 0.83--0.89 and does not restate sign
agreement or coverage. Recorded so this is not reopened. Consequence to keep
in view, that is the number the meta-review turned into ``explain 90\% of the
variance of the observations'', and the previous submission did carry ``at
least 91\% of logit differences'' at old L23--24. The two headline metrics
are stated in this document instead, in the j4yY abstract bullet, so a
resubmission reader meets them early even though the abstract does not carry
them. Do not raise the meta-review's phrasing anywhere.
\item[$\boxtimes$] Label-generation pipeline. Done 2026-08-04. Appendix~L,
\emph{How labels are produced}, states that a language model drafts each
description from the feature's highest-scoring unembedding rows and the
authors verify every claim-bearing label against those same rows, and it
frames the audit as the measurement of that verification's success rate.
pZ3j's W2 premise that labelling is ``all done by humans'' no longer stands
unanswered, and checklist item E1 is discharged. The W2 bullet in the
point-by-point opens with the correction.
\item[$\boxtimes$] Contrast banks and blinded-audit protocol in Appendix~F.
Verified 2026-08-04: the artifact-availability paragraph names both, in both
the preprint and the anonymous branch, and the URLs are correctly guarded by
\texttt{\textbackslash ifpreprintmode}. The Section~3 artifacts bullet now
names them.
\item[$\boxtimes$] $k=32$. Verified 2026-08-04: Appendix~D.3
(\emph{Lower Sparsity Budgets}) already reports $k=32$ \emph{and} $k=64$ on
both Qwen3.5-2B and Qwen3.5-0.8B. Nothing to run. The earlier belief that
only the $k=64$ Qwen3.5-2B cells existed was wrong, and this document had
been understating the experiment; both it and the map table now say two
models at two budgets. Note $k=32$ carries 1 seed and rowEV is withheld where
the dead rate exceeds the selection threshold.
\item[$\boxtimes$] Appendix~I.4 and Appendix~L retain what this document no
longer restates. Verified 2026-08-04: I.4 carries the cross-recipe value
(0.444, below row-kNN) and L records the earlier three-way label scheme.
\item[$\boxtimes$] Add the $\mu$ justification. Done in Appendix~A.1,
paragraph \emph{Why the offset stays outside the dictionary}.
\item[$\boxtimes$] Explain why the oracle's KL of 0.004 is not an
information-matched comparison. Done in Appendix~N.
\item[$\boxtimes$] Reading-rule pass rate. Resolved 2026-08-04:
\texttt{results\_and\_analysis.tex} defines coverage as the fraction of
scores the reading rule accepts, so the thread's 69--89\% and the paper's
0.72--0.81 are the same quantity on different banks. This document quotes
only the paper's, so no action.
\item \textbf{Consider.} Appendix~I.4 is titled \emph{Cross-Recipe
Disclosure}. ``Disclosure'' frames a scoping result as a confession in the
table of contents. \emph{Cross-Recipe Comparison} says the same thing.
\end{itemize}

\textbf{This document}
\begin{itemize}[label=$\square$]
\item[$\boxtimes$] Placeholders. Resolved 2026-08-04 by dropping line
references entirely. Locations are now section and appendix numbers, which
are stable, so this document no longer depends on a frozen \texttt{main.pdf}
and \texttt{main.tex} can be recompiled freely. Zero
\texttt{\textbackslash ph\{\}} remain. Do not reintroduce line numbers.
\item[$\boxtimes$] Old line numbers. Verified 2026-08-04 against
\texttt{../10103\_Sparse\_Readout\_Prism\_A\_S.pdf}, the previous submission.
L24 ``on at least 91\% of logit differences'', L84 ``on hidden-state
activations. Once a readout score'', L196 ``is positive when the readout
favors token A over'', L302 ``scores survive the factorization. Sign
agree-''. All four match the reviews' citations.
\item Set \texttt{\textbackslash draftnotesfalse}. This removes every red
note and this checklist.
\item Then verify the posting PDF's text layer is clean. Extract it and grep
for \texttt{TODO}, \texttt{NOT IN PAPER}, \texttt{DECISION} and
\texttt{meta-reviewer}. Upload the compiled PDF only, never this source.
\item[$\boxtimes$] Anonymity pass on this document. Verified 2026-08-04: no
author names, affiliations, repository URLs or author-list change
information.
\item \textbf{OPEN.} Mark \texttt{RESUBMISSION\_CHANGE\_SUMMARY.md} as
superseded or delete it. Confirmed still present 2026-08-04. It covers the
same ground with stale appendix letters and table numbers.
\item \textbf{OPEN, page limit has no slack.} Verified 2026-08-04 from a
fresh build: Conclusion starts on p.~8 and Limitations on p.~9, so the body
is compliant with nothing to spare. Any addition to \S1--\S6, including the
abstract fix above, must be paid for by a cut of similar length.
\item Upload to the Explanation of Revisions field. Optionally prepend to the
main submission PDF.
\end{itemize}

\textbf{Submission metadata, not part of this PDF}
\begin{itemize}[label=$\square$]
\item Set the two Reassignment Request fields. Both currently read ``This is
not a resubmission.''
\item Update \emph{Languages Studied}, currently English only, since \S4 now
covers English, Chinese and German.
\item Re-answer checklist C1 and C4 as Yes, pointing at Appendix~F.
\item Confirm the preprint decision. The previous submission declared an
intent to post a non-anonymous preprint during review, so check this against
the resubmission's anonymity requirements before posting.
\end{itemize}
}
\fi

\end{document}
